% Source-local optional insertion for the first use of the
% non-homogeneous-host assumption in Lemma 6.6 / Proposition 6.17.
% Independent author review only. The circulated manuscript is frozen.

If $H$ were a disjoint union of cliques, vertex-transitivity
would force all its components to have the same order.
Such a graph is homogeneous: any partial automorphism induces
an injection between the components meeting its domain
and those meeting its image. This extends to a permutation
of the components, and the prescribed partial bijections
extend inside them since the components have equal order.
Consequently the non-homogeneous graph $H$ is not a
disjoint union of cliques.

% Rationale: the source assumes non-homogeneous vertex-transitive H,
% while the strongest Lean version of the generalized-cover and
% unequal-clique theorems assumes non-cluster H. The displayed
% elementary implication supplies the exact logical bridge.
% This is a paper proof, not yet an exported generic Lean theorem.
